\section{问题重述}

\subsection{研究背景}

舰船在复杂海空威胁环境中需要在有限的反应时间内完成侦测、释放和机动。烟幕弹的效果并不是一个静态几何覆盖问题：目标与来袭威胁都在运动，烟幕云团随时间膨胀、保持和衰减，释放平台还受到航迹长度、反应延迟和作用半径的限制。因此，一个仅在单个时刻判断“是否遮蔽”的模型，不能直接回答连续时间内的遮蔽持续性和最坏暴露问题。

本文将题目中的决策链拆解为四个层次。第一层研究单枚烟幕弹的连续起爆时刻和中心时刻；第二层研究同一架 UAV 携带多枚烟幕弹时的事件组合与联合覆盖；第三层研究多架 UAV 分担烟幕弹后的航迹可行性和协同覆盖；第四层把前述已认证任务包交给一个带有揭示时刻的多威胁调度器。

\subsection{四个子问题}

\begin{enumerate}
  \item \textbf{单弹连续优化（Q1）}：在给定威胁场景、舰船运动和反应延迟下，寻找释放/起爆相关的连续决策，使遮蔽覆盖更长、最大连续暴露更短，同时保留航迹约束。需要比较四类候选算法，并使用独立验证器复核。
  \item \textbf{单机多弹联合覆盖（Q2）}：同一架 UAV 依次释放至多三枚烟幕弹。需要把事件时刻限制到有限候选集，再进行局部精修，使用保守的连续时间覆盖证书区分已覆盖、已暴露和未解析区间。
  \item \textbf{多机协同（Q3）}：固定三架 UAV，每架恰好一枚烟幕弹。需要检查每架 UAV 的作用半径、任意两架之间的路径冲突，并以带 $\varepsilon$ 阈值的词典序规则选择联合方案。
  \item \textbf{多威胁因果调度（Q4）}：从 Q1--Q3 已认证的任务包出发，在任务揭示过程中比较因果 rolling、因果 greedy 与离线 hindsight 上界；未揭示或未认证的任务不能提前消耗容量。
\end{enumerate}

本文的边界是“可审计的有限搜索与验证”。论文不把有限预算算法的最好候选写成连续全局最优，也不把离线 hindsight 上界写成在线策略；这些限制直接由结果文件的 \texttt{global\_optimum\_claim=false}、Q4 的 \texttt{causal} 字段和独立验证代码体现。

\subsection{研究对象与输出}

研究对象为四个正式场景中的代表性前向威胁矩阵行，以及 Q1 生成的 $4\times4$ 正式矩阵。输出包括候选决策、覆盖下界/上界、总暴露与最大连续暴露、最小安全裕度、航迹距离、路径冲突证书、资源容量下的调度值和局部灵敏度结果。所有论文数值均来自 \path{results/q1_rebuild/}、\path{results/q2_rebuild/}、\path{results/q3_rebuild/}、\path{results/q4_rebuild/} 和 \path{results/sensitivity_rebuild/}，证据等级与允许措辞见论文审计表。
